% RPLU v2 / JET Pipeline Paper — IEEE Micro / FPGA / ASPLOS target
% A Hardware Jet Algebra Coprocessor over a Split M31 Biquadratic Algebra:
% Higher-Order Differentiation without Floating-Point
\documentclass[conference]{IEEEtran}
\usepackage{cite}
\usepackage{amsmath,amssymb,amsfonts}
\usepackage{graphicx}
\usepackage{booktabs}
\usepackage{multirow}
\usepackage[colorlinks=true]{hyperref}

\begin{document}

\title{A Hardware Jet Algebra Coprocessor over a Split M31 Biquadratic Algebra:
Deterministic Higher-Order Differentiation without Floating-Point}

\author{
\IEEEauthorblockN{John G.}
\IEEEauthorblockA{SPU-13 Project\\
github.com/pinnacletechnologysolutionsltd/SPU}
}

\maketitle

\begin{abstract}
Higher-order differentiation in embedded and edge computing is dominated by
finite-difference approximations and IEEE-754 floating-point accumulators,
both of which introduce rounding drift, non-deterministic convergence, and
variable-latency subnormals. This paper presents a hardware jet algebra
coprocessor that evaluates exact Cauchy products over the truncated
polynomial ring $A_{31}[\varepsilon]/(\varepsilon^3)$ using the
Mersenne prime $p = 2^{31}-1$ (M31), where
$A_{31}=\mathbb{F}_p[u,v]/(u^2-3,v^2-5)$ is the split biquadratic algebra
implemented by the RTL.  A parameterized jet multiply-accumulate
(MAC) pipeline computes $R_k = \sum_{i=0}^{k} J_i \cdot K_{k-i}$ for $k \in
\{0,1,2\}$ in 12 cycles (time-division multiplexed) or 3 cycles (fully
unrolled), with bit-identical replay across simulation runs. A conjugate reduction tower and shadow chain invert unit elements in
$\sim$88 deterministic cycles total ($\sim$76 for the tower, 12 for the chain). Issue control is deterministic and blocking:
multi-cycle tower operations gate issue through busy/done clock-enable
handshakes, while single-cycle arithmetic avoids read-after-write stalls
through a write-forwarding register-file bypass.

The leaf RTL for the 4-stage pipeline---jet MAC, jet inverse, [4/4] Pad\'{e} rational
approximant, and Kohonen SOM classifier---is under active integration, with
focused Verilog benches and Python/C++$\,$17 oracles covering the arithmetic
and control primitives. Two hardware targets are presented:
a Xilinx Artix-7 100T implementation (240 DSP48E1, 63,400 6-input LUTs)
currently routed at a 4.11 MHz bring-up frequency, and a Gowin GW5A-25A
silicon-proven prototype (8,256 LUT4s) demonstrating subsystem feasibility
at 67\% device utilization. The architecture eliminates
floating-point non-determinism, square roots, and runtime dynamic
address resolution from the topological classification hot path.
All RTL, software oracles, and documentation are open-source (CC0 1.0).
\end{abstract}

% ── SECTION I: Mathematical Foundations ────────────────────────────

\section{Mathematical Foundations}
\label{sec:math}

\subsection{The Split Biquadratic Algebra $A_{31}$ over M31}

The arithmetic foundation is the Mersenne prime $p = 2^{31} - 1 =
2147483647$ (M31), chosen for its hardware-friendly reduction: $x \bmod p
= x_{\text{lo}} + x_{\text{hi}}$ with a single conditional subtract for
operands up to 72 bits. The current RTL uses the algebra
\begin{equation}
A_{31}=\mathbb{F}_p[u,v]/(u^2-3,\;v^2-5)
\end{equation}
with basis $[1,u,v,uv]$, historically written in the source as
$[1,\sqrt{3},\sqrt{5},\sqrt{15}]$.

This is intentionally described as a finite commutative algebra, not as
$\mathbb{F}_{p^4}$. Euler's criterion gives $3^{(p-1)/2}\equiv -1$ and
$5^{(p-1)/2}\equiv -1$, but $15^{(p-1)/2}\equiv 1$; concretely,
$1393679181^2\equiv 15\pmod p$. Thus $A_{31}$ contains structural
zero-divisors. The hardware treats zero-norm elements as non-units and routes
them to the singular absorber rather than silently using a false inverse.

An element $Z \in A_{31}$ is a 4-tuple $(z_0, z_1, z_2, z_3) \in
[0, p-1]^4$ representing $Z = z_0 + z_1u + z_2v + z_3uv$.

Multiplication follows the cross-product rules in Table~\ref{tab:mult-rules}.

\begin{table}[ht]
\centering
	\caption{$A_{31}$ multiplication rules.}
\label{tab:mult-rules}
\begin{tabular}{c|cccc}
$\times$ & $1$ & $\sqrt{3}$ & $\sqrt{5}$ & $\sqrt{15}$ \\
\midrule
$1$      & $1$    & $\sqrt{3}$ & $\sqrt{5}$ & $\sqrt{15}$ \\
$\sqrt{3}$  & $\sqrt{3}$ & $3$       & $\sqrt{15}$  & $3\sqrt{5}$ \\
$\sqrt{5}$  & $\sqrt{5}$ & $\sqrt{15}$  & $5$        & $5\sqrt{3}$ \\
$\sqrt{15}$ & $\sqrt{15}$ & $3\sqrt{5}$ & $5\sqrt{3}$ & $15$
\end{tabular}
\end{table}

\subsection{Rational Quadrance: Purging Square Roots from Silicon}

In a standard Euclidean Self-Organizing Map (SOM), the distance between a
feature vector $\mathbf{x}$ and a node weight vector $\mathbf{w}_i$ is:

\begin{equation}
d(\mathbf{x}, \mathbf{w}_i) = \sqrt{\sum_j (x_j - w_{ij})^2}
\end{equation}

The square root is (a) non-deterministic in floating-point, (b)
power-hungry, and (c) irrelevant for \emph{ordering}, since $\sqrt{a} <
\sqrt{b} \iff a < b$. The RPLU replaces this with the rational quadrance
metric:

\begin{equation}
Q(\mathbf{x}, \mathbf{w}_i) = \sum_j r_j \cdot (x_j - w_{ij})^2
\end{equation}

	where $r_j \in A_{31}$ are rational feature weights and all
operations are exact modular algebra elements. No square root, no floating-point, no
approximation. The winner-take-all (WTA) comparator tree operates on exact
integer comparisons over the quadrance values.

\subsection{Nguyen Tolerance Rough Sets and Information Granules}

Hung Son Nguyen's tolerance rough sets~\cite{nguyen1999} provide the
theoretical framework for cluster ambiguity. In traditional SOM, a
``tie'' between two BMU candidates requires domain-specific heuristics
or random arbitration. The RPLU computes an explicit Confidence Gap:

\begin{equation}
\text{Gap} = Q_{\text{second-best}} - Q_{\text{best}}
\end{equation}

When $\text{Gap} = 0$ (ambiguous), the hardware classifies the point as
belonging to a \emph{boundary region}—a hardware-driven Information
Granule that defers classification rather than making a brittle
decision. This is a first-class architectural primitive, not a post-hoc
software correction.

% ── SECTION II: Jet Algebra Microarchitecture ───────────────────────

\section{Jet Algebra Microarchitecture}
\label{sec:jet}

	\subsection{The Jet Ring $A_{31}[\varepsilon]/(\varepsilon^3)$}

The core architectural primitive is not the SOM or the Pad\'{e} approximant
	but the \emph{jet}---a truncated polynomial expansion over the local ring
	$A_{31}[\varepsilon]/(\varepsilon^3)$, where
$\varepsilon^3 \equiv 0$. A jet $J \in \mathbb{A}_{\text{SPU}}$ is a
3-tuple over $A_{31}$:

\begin{equation}
J = (j_0, j_1, j_2) \quad\text{where}\quad
	j_k \in A_{31},\; k \in \{0, 1, 2\}
\end{equation}

representing the truncated series $j_0 + j_1\varepsilon + j_2\varepsilon^2$.
The coefficient $j_0$ encodes the \emph{position} (base value), $j_1$ the
\emph{velocity} (first derivative), and $j_2$ the \emph{acceleration}
(second derivative) of a quantity propagating through the pipeline. All
three coefficients share the same $A_{31}$ modular arithmetic
described in Section~\ref{sec:math}.

The truncation is made before evaluation. RPLU does not evaluate an infinite
Taylor series and then discard high-order terms; it changes the algebraic
object to the quotient ring $A_{31}[\varepsilon]/(\varepsilon^{N+1})$ first.
Terms of degree greater than $N$ are identically zero in that ring, so the
Cauchy products are exact finite operations. Transcendental or Pad\'{e}
blueprints enter as finite coefficient sets, and runtime failure is limited
to explicit unit-inversion traps in $A_{31}$.

\subsection{Jet Multiply-Accumulate (Jet MAC)}

Given two jets $J$ and $K$, their Cauchy product $R = J \cdot K$ is:

\begin{equation}
R_k = \sum_{i=0}^{k} J_i \cdot K_{k-i} \quad\text{for } k \in \{0, 1, 2\}
\end{equation}

Expanding:
\begin{align}
R_0 &= J_0 \cdot K_0 \\
R_1 &= J_0 \cdot K_1 + J_1 \cdot K_0 \\
R_2 &= J_0 \cdot K_2 + J_1 \cdot K_1 + J_2 \cdot K_0
\end{align}

	Each $\cdot$ is a full $A_{31}$ multiplication requiring 16
parallel 32$\times$32 integer products (4 ring components $\times$ 4 basis
cross-terms). The jet MAC (\texttt{spu13\_jet\_mac.v}) is parameterized by
order $N$ (default $N=2$ for $\varepsilon^3$ truncation) and implements
these six multiplications as a time-division multiplexed (TDM) sequence
through a shared M31 multiplier:

\begin{equation}
\text{Latency}_{\text{MAC}} = 2 \cdot \frac{(N+1)(N+2)}{2} = 12\ \text{cycles}
\end{equation}

For $N=2$, the 6 sequential multiplies each take 2 cycles (one cycle
multiply, one cycle Mersenne reduction + accumulate), giving 12-cycle total
latency. A fully unrolled variant replicates the multiplier 6$\times$ for
single-cycle Cauchy completion.

\subsection{Jet Inverse (Shadow Chain)}

	The jet inverse $J^{-1}$ of a jet whose base coefficient $j_0$ is a unit is computed
by solving $J \cdot J^{-1} = (1, 0, 0)$ through a shadow chain of
sequential multiplies and additions. The inverter
(\texttt{spu13\_jet\_inv.v}) decomposes inversion into:

\begin{align}
m_0 &= j_0^{-1} \quad\text{(via conjugate reduction tower)} \\
	m_1 &= -j_1 \cdot m_0^2 \\
	m_2 &= j_1^2 \cdot m_0^3 - j_2 \cdot m_0^2
\end{align}

where $j_0^{-1}$ is computed by the 76-cycle conjugate reduction tower
(Section~\ref{sec:tower}) and $m_0, m_1, m_2$ are the inverse jet
coefficients. The shadow chain adds $\sim$12 cycles after tower completion,
for a total inversion latency of $\sim$88 cycles.

	If $j_0$ is a non-unit (zero norm), the hardware asserts \texttt{err\_zero\_divisor}
and flags the result as a geometric singularity---no hang, no silent corruption.

\subsection{The PHSLK Instruction}

The Wheeler-Feynman phase-lock instruction (PHSLK, opcode 0x42) is the
top-level architectural expression of the jet pipeline. It takes two jet
operands---an Offer jet $O$ and a Confirmation jet $C$---and computes
the coherence predicate:

\begin{equation}
\text{Coherent} \iff O \cdot C \equiv (1, 0, 0) \in \mathbb{A}_{\text{SPU}}
\end{equation}

Decomposed into three lane checks:
\begin{align}
\text{Lane 0 (base)} &: o_0 \cdot c_0 \equiv 1 \\
\text{Lane 1 (velocity)} &: o_0 \cdot c_1 + o_1 \cdot c_0 \equiv 0 \\
\text{Lane 2 (acceleration)} &: o_0 \cdot c_2 + o_1 \cdot c_1 + o_2 \cdot c_0 \equiv 0
\end{align}

All three lanes are computed through the shared jet MAC. When all three
pass, \texttt{FLAGS.C = 1} (locked) and the pipeline may proceed. This
is a single hardware instruction---no branches, no microcode, no
floating-point comparison.

\subsection{From Jets to the RPLU Pipeline}

The RPLU v2 pipeline (Section~\ref{sec:microarch}) is an application of
this jet algebra substrate. The SOM/BMU classifier produces a spatial
index; the BTU routes that index into an $A_{31}$
element; the [4/4] Pad\'{e} approximant evaluates a rational function;
and the conjugate reduction tower provides the $j_0^{-1}$ required by the
jet inverse. All four stages operate on the same
$A_{31}[\varepsilon]/(\varepsilon^3)$ jet representation where the jet path is
enabled, with the same fixed-latency, bit-identical arithmetic.

% ── SECTION III: Pipeline Microarchitecture ─────────────────────────

\section{Pipeline Microarchitecture}
\label{sec:microarch}

The RPLU v2 is a 4-stage pipeline:

\begin{equation}
\Phi_1\text{: SOM BMU} \rightarrow
\Phi_2\text{: BTU router} \rightarrow
\Phi_3\text{: Pad\'{e} + Inv} \rightarrow
\Phi_4\text{: Output latch}
\end{equation}

\subsection{Stage $\Phi_1$ — Ingestion Layer: The Bus Transmutation Unit}

The BTU (Barycentric Transmutation Unit) is a Direct Combinational Demultiplexer
Grid that executes $O(1)$ spatial-to-bus mapping. A 64-bit activation
vector arrives from the upstream SOM classification stage. Each bit
represents a neutron in the Kohonen mesh. The BTU performs a 64-to-6
priority encoding that selects the lowest-index active neutron as the
current dispatch target.

The spatial address space is hierarchically organized: 4 quadrants, each
with 4 clusters, each with 4 neurons. The 6-bit output encodes (2-bit
quadrant, 2-bit cluster, 2-bit neuron). This mapping feeds directly into
the 4-lane BRAM lookup that maps the spatial index into an
$A_{31}$ element.

\subsection{The Control Layer: Multi-Hot Collision Resolution}

In physical wave-interference topologies, multiple saddle points may be
active simultaneously---a \emph{multi-hot} activation pattern. The
collision resolver (\texttt{spu\_btu\_collision\_resolver.v}) handles this
deterministically:

\begin{enumerate}
\item \textbf{Priority encoder:} 64 $\rightarrow$ 6 combinational
selection, lowest-index active neuron wins.
\item \textbf{Backlog queue:} Remaining active bits are queued in a
shift-register backlog. The pipeline is \emph{bubble-stalled} by asserting
\texttt{pipeline\_stall} upstream.
\item \textbf{Serial dispatch:} Each subsequent cycle dequeues the next
active neutron. Latency is $O(n)$ where $n$ is the number of simultaneous
activations.
\end{enumerate}

Unlike a round-robin arbiter that requires state tracking, this priority
encoder is purely combinational and stateless---the ordering is
\emph{topological} (lowest spatial index first), not temporal. This
guarantees identical dispatch order for identical activation patterns
across simulation runs; post-place timing remains a silicon-integration task.

\subsection{Stage $\Phi_3$ — The Conjugate Reduction Tower}
\label{sec:tower}

Unit inversion in $A_{31}$ avoids brute-force exponentiation by exploiting the
nested quadratic structure:
\begin{enumerate}
\item \textbf{Relative norm to $\mathbb{F}_p[u]/(u^2-3)$:} Write $Z = A +
Bv$ where $A, B \in \mathbb{F}_p[u]/(u^2-3)$. Compute $D = A^2 -
5B^2$.
\item \textbf{Norm to $\mathbb{F}_p$:} $N = d_0^2 - 3d_1^2 \in \mathbb{F}_p$.
\item \textbf{Zero-norm trap:} If $N \equiv 0$, assert FLAGS.V and bypass
the Fermat chain (non-unit / zero-divisor manifold).
\item \textbf{Fermat reconstruction:} $N^{-1} = N^{p-2} \bmod p$ via
hardwired 30-step square-and-multiply.
\item \textbf{Recombination:} $Z^{-1} = (A - Bv) \cdot (d_0 -
d_1(-u)) \cdot N^{-1}$.
\end{enumerate}

Total latency: $\sim 76$ cycles, completely deterministic, zero division
operations. Inversion test vectors are shown in Table~\ref{tab:fp4-inverter}.

\paragraph{Montgomery batch inversion.}
When the Pad\'{e} evaluator must invert $k$ denominators---as in a
multi-lane manifold sweep or trajectory propagation---the per-element
tower cost of $k \cdot 76$ cycles is dominated by the serial Fermat-chain
latency. Montgomery's product-accumulate/unwind trick~\cite{montgomery1987}
collapses $k$ inversions to \textbf{1 tower + $3(k-1)$ multiplies} by
computing a single prefix-product chain over all denominators, inverting
the accumulated product once, then unwinding the individual inverses via
the recurrence $d_i^{-1} = (d_0\cdots d_{k-1})^{-1} \cdot (d_0\cdots
d_{i-1}) \cdot (d_{i+1}\cdots d_{k-1})$. The zero-divisor check is
deferred: the norm is multiplicative into $\mathbb{F}_p$, so the
accumulated product's norm is zero iff at least one factor is singular.
One tower-norm stage covers the whole batch; on a hit, a per-element norm
probe ($2$ multiplies each, stages A+B only, no Fermat) isolates the
offending lanes, which then receive per-lane FLAGS.V semantics while
unit lanes emerge bit-exact. The crossover at which batch inversion
breaks even is $k = 2$; the asymptote is $2.86\times$ speedup with a
$+33\%$ MAC volume ceiling. Table~\ref{tab:batch-inv} quantifies the
trade-off on representative RPLU2 workload mixes.

\subsection{The Execution Layer: Deterministic Issue Control}

The 4R2W multi-port register file (\texttt{spu13\_multi\_port\_regfile.v})
supports single-cycle integer arithmetic (QADD, QSUB, QMUL) beside the
76-cycle Conjugate Inversion Tower.

\textbf{Key design choice:} Single-cycle ops may issue every cycle without
stalling---they write through the register file's combinational
write-forwarding bypass. Tower ops (PHSLK, INVJ, Pad\'{e} eval) are
\emph{blocking}: issue is gated by the tower's busy/done clock-enable
handshake until writeback completes. Since the tower latency is fixed,
all hazard resolution is combinational---no dynamic arbitration, reorder
buffers, or register scoreboarding are required. A prototype scoreboard
module for concurrent issue during tower runs exists in the repository
but is not instantiated in any verified build; scoreboarded overlap of
independent work with tower latency is future work
(Section~\ref{sec:future}).

The write-forwarding bypass means a result written on cycle $t$ is
available to the instruction issuing on cycle $t+1$ via a mux-select path
rather than a register-file re-read. This eliminates the read-after-write
stall that would otherwise add a cycle to every dependent instruction pair.

\subsection{The [4/4] Pad\'{e} Rational Approximant}

The Pad\'{e} stage evaluates:

\begin{equation}
R(x) = \frac{p_4 x^4 + p_3 x^3 + p_2 x^2 + p_1 x + p_0}
            {q_4 x^4 + q_3 x^3 + q_2 x^2 + q_1 x + q_0}
\end{equation}

via Horner's method: $(\cdots(p_4 \cdot x + p_3) \cdot x + \cdots + p_0)$.
The denominator is inverted via the Conjugate Reduction Tower. Numerator
and denominator coefficients are stored in 2$\times$5$\times$32-bit BRAM
banks, loaded at boot via the config port.

The singular absorber condition---when the denominator norm approaches
zero---is handled by the same zero-norm trap in the tower. The hardware
asserts FLAGS.V and drops \texttt{done} cleanly rather than hanging on a
near-singular division.

% ── SECTION III: Verification ──────────────────────────────────────

\section{Verification Methodology}
\label{sec:verification}

\subsection{Co-Design Verification Workflow}

The verification strategy uses three independent implementations of the
same arithmetic:

\begin{enumerate}
\item \textbf{Python oracles} (\texttt{rational\_som.py} + \texttt{spu13\_arch\_sim.py}):
High-level behavioral models covering SOM BMU classification, stable
tie-breaking, confidence gap calculation, and full jet algebra execution.
24 test checks.
\item \textbf{C++ oracles} (\texttt{spu\_rational\_som.h} + \texttt{spu13\_arch.h}):
Cycle-accurate parity with Python oracles. 51 test checks.
\item \textbf{RTL simulation} (Icarus Verilog): focused benches covering
jet MAC, jet INV, PHSLK coherence, $A_{31}$ multiplication,
conjugate reduction tower, BTU collision, SOM BMU, and Pad\'{e} evaluation.
\end{enumerate}

All three layers use identical 32-bit $A_{31}$ vector encodings.
The Python and C++ oracles directly mirror the clock-edge execution
semantics of the Verilog testbenches---e.g., the 2-stage multiplier
pipeline latency is modeled as a function call with a cycle counter,
not as a single arithmetic operation.

\subsection{Jet Pipeline Verification}

The jet MAC testbench (\texttt{spu13\_jet\_mac\_tb}) validates 3 scenarios:
unit propagation ($J = (1,0,0)$, $K = (1,0,0)$), zero-divisor guard
($j_0 = 0$), and full Cauchy product with random coefficients. The jet
inverse testbench (\texttt{spu13\_jet\_inv\_tb}) validates 6 scenarios:
scalar identity $J^{-1}$ where $J = (5,0,0)$, inversion with velocity
$J = (5,3,2)$, zero-divisor trap $J = (0,1,0)$, post-trap recovery,
and random self-consistency $J \cdot J^{-1} = (1,0,0)$. The PHSLK core
testbench (\texttt{spu13\_phslk\_core\_tb}) validates 6 coherence
predicate scenarios: matching base, mismatched base, zero-divisor,
post-trap valid lock, velocity cancellation, and acceleration
cancellation. The full-stack PHSLK testbench (\texttt{spu\_full\_stack\_tb})
validates 13 additional scenarios including concurrent LDO/LDC/PHSLK
sequences, write-forwarding bypass, and RAW hazard resolution.

The current jet MAC leaf bench exercises exact $\varepsilon^3$ truncation.
The jet inverse and PHSLK benches are being expanded to assert full
self-consistency ($J\cdot J^{-1}=(1,0,0)$) rather than only nonzero outputs.

\subsection{Cross-Validation Results}

Table~\ref{tab:verification} summarizes the verification layers.
These tables summarize the current oracle and leaf-RTL coverage. They are not
yet a substitute for integrated pipeline elaboration, timing closure, or
silicon measurements.

% ── RPLU v2 Publication Data — Auto-generated ──

% Table 1: A31 Conjugate Reduction Tower — Unit Inversion Vectors
\begin{table}[ht]
\centering
\caption{$A_{31}$ conjugate reduction tower unit-inversion vectors generated
by the publication oracle. The corresponding inverter bench is included in
the full RTL gate. Input $Z = (z_0, z_1, z_2, z_3) \in A_{31}$ over $p = 2^{31}-1$
(M31). Non-units, including nonzero zero-divisors, trap via FLAGS.V.
$\sim$76 cycle deterministic latency.}
\label{tab:fp4-inverter}
\resizebox{\columnwidth}{!}{%
\begin{tabular}{llll}
\toprule
Test Case & Input $Z$ & Oracle $Z^{-1}$ & Match \\
\midrule
  Identity & (00000001,00000000,00000000,00000000) & (00000001,00000000,00000000,00000000) & $\checkmark$ \\
  Scalar inv(2) & (00000002,00000000,00000000,00000000) & (40000000,00000000,00000000,00000000) & $\checkmark$ \\
  Pure $\sqrt{3}$ inv & (00000000,00000001,00000000,00000000) & (00000000,55555555,00000000,00000000) & $\checkmark$ \\
  Pure $\sqrt{5}$ inv & (00000000,00000000,00000001,00000000) & (00000000,00000000,33333333,00000000) & $\checkmark$ \\
  Zero-norm & (00000000,00000000,00000000,00000000) & SINGULARITY & FLAGS.V \\
  Nonzero zero-divisor & (2CEE24B2,00000000,00000000,00000001) & SINGULARITY & FLAGS.V \\
  Random self-consistency & (00003039,00010932,00002B67,000056CE) & (7122397B,1212AF65,4D457009,1D8AB3FE) & N/A \\
\bottomrule
\end{tabular}
}%
\end{table}



% Table 2: A31 Multiplication — M31 Mersenne Reduction
\begin{table}[ht]
\centering
\caption{$A_{31}$ split-biquadratic multiplication over M31. Sixteen logical
$32\times32$ products feed a 2-stage pipeline with Mersenne reduction via
72-bit chunk splitting. Physical DSP usage is target- and synthesis-dependent.}
\label{tab:m31-mult}
\resizebox{\columnwidth}{!}{%
\begin{tabular}{lllll}
\toprule
Test Case & $A$ & $B$ & $A \cdot B \bmod p$ & Match \\
\midrule
  Identity & (00000001,00000000,00000000,00000000) & (00000001,00000000,00000000,00000000) & (00000001,00000000,00000000,00000000) & $\checkmark$ \\
  Scalar 5$\times$3 & (00000005,00000000,00000000,00000000) & (00000003,00000000,00000000,00000000) & (0000000F,00000000,00000000,00000000) & $\checkmark$ \\
  $\sqrt{3}$ $\times$ $\sqrt{3}$ = 3 & (00000000,00000001,00000000,00000000) & (00000000,00000001,00000000,00000000) & (00000003,00000000,00000000,00000000) & $\checkmark$ \\
  $\sqrt{5}$ $\times$ $\sqrt{5}$ = 5 & (00000000,00000000,00000001,00000000) & (00000000,00000000,00000001,00000000) & (00000005,00000000,00000000,00000000) & $\checkmark$ \\
  $\sqrt{15}$ $\times$ $\sqrt{15}$ = 15 & (00000000,00000000,00000000,00000001) & (00000000,00000000,00000000,00000001) & (0000000F,00000000,00000000,00000000) & $\checkmark$ \\
  $\sqrt{3}$ $\times$ $\sqrt{5}$ = $\sqrt{15}$ & (00000000,00000001,00000000,00000000) & (00000000,00000000,00000001,00000000) & (00000000,00000000,00000000,00000001) & $\checkmark$ \\
  M31 edge: (P-1)$\times$2 & (7FFFFFFE,00000000,00000000,00000000) & (00000002,00000000,00000000,00000000) & (7FFFFFFD,00000000,00000000,00000000) & $\checkmark$ \\
  Mixed: (10,2,0,4)$\times$(5,0,1,2) & (0000000A,00000002,00000000,00000004) & (00000005,00000000,00000001,00000002) & (000000AA,0000001E,00000016,0000002A) & $\checkmark$ \\
  Zero multiply & (00000000,00000000,00000000,00000000) & (0000007B,000001C8,00000315,00018AF8) & (00000000,00000000,00000000,00000000) & $\checkmark$ \\
\bottomrule
\end{tabular}
}%
\end{table}



% Table 3: M31 Scalar Inverter — BEEA Test Vectors
\begin{table}[ht]
\centering
\caption{M31 scalar modular inversion via binary extended Euclidean algorithm:
zero divisions, using shifts and conditional $p$ addition.}
\label{tab:m31-inv}
\resizebox{\columnwidth}{!}{%
\begin{tabular}{llll}
\toprule
Test Case & $x$ & $x^{-1} \bmod p$ & Match \\
\midrule
  inv(1) & 1 & 00000001 & $\checkmark$ \\
  inv(2) & 2 & 40000000 & $\checkmark$ \\
  inv(P-1) & 2147483646 & 7FFFFFFE & $\checkmark$ \\
  inv(3) & 3 & 55555555 & $\checkmark$ \\
  inv(1234567) & 1234567 & 37B9E691 & $\checkmark$ \\
  inv(65537) & 65537 & 0000FFFF & $\checkmark$ \\
\bottomrule
\end{tabular}
}%
\end{table}



% Table 4: Singular Absorber — Zero-Norm Exception Path
\begin{table}[ht]
\centering
\caption{Publication-oracle checks for FLAGS.V behavior on zero-norm
singularities and clean re-arm after exception. The singular-absorber RTL
bench is included in the full release gate.}
\label{tab:singular-absorber}
\resizebox{\columnwidth}{!}{%
\begin{tabular}{lllll}
\toprule
Scenario & Input $Z$ & Result & FLAGS.V & Match \\
\midrule
  Valid inv(5) & (00000005,00000000,00000000,00000000) & (33333333,00000000,00000000,00000000) & 0 & $\checkmark$ \\
  Zero $\to$ singularity & (00000000,00000000,00000000,00000000) & TRAP & 1 & $\checkmark$ \\
  Unity (false positive guard) & (00000001,00000000,00000000,00000000) & (00000001,00000000,00000000,00000000) & 0 & $\checkmark$ \\
  Re-arm zero & (00000000,00000000,00000000,00000000) & TRAP & 1 & $\checkmark$ \\
  Post-exception inv(3) & (00000003,00000000,00000000,00000000) & (55555555,00000000,00000000,00000000) & 0 & $\checkmark$ \\
\bottomrule
\end{tabular}
}%
\end{table}



% Table 5: BTU Collision Resolver — Multi-Hot Wave Dispatch
\begin{table}[ht]
\centering
\caption{BTU collision resolver (64$\to$6 priority encoder + backlog queue)
serializing multi-hot wave interference into deterministic dispatch.
Fixed O(n) latency where n = number of simultaneous activations.}
\label{tab:btu-collision}
\resizebox{\columnwidth}{!}{%
\begin{tabular}{llll}
\toprule
Scenario & Activation bits & Cycles & Dispatch sequence \\
\midrule
  Single node (bit 3) & 3 & 1 & k=3 \\
  Two nodes (bits 7,15) & 7, 15 & 2 & k=7 (stall) $\to$ k=15 \\
  Three nodes (bits 0,33,63) & 0, 33, 63 & 3 & k=0 (stall) $\to$ k=33 (stall) $\to$ k=63 \\
  Zero activation &  & 0 & idle, no stall \\
\bottomrule
\end{tabular}
}%
\end{table}



% Table 6: SOM BMU Classification — Rational Quadrance Results
\begin{table}[ht]
\centering
\caption{BMU classification using weighted rational quadrance metric
(no square roots, no floating-point). 7-node parallel array with
combinational winner-take-all tree. 3-stage per-node pipeline:
subtract$\to$square$\to$accumulate.}
\label{tab:som-bmu}
\resizebox{\columnwidth}{!}{%
\begin{tabular}{llllll}
\toprule
Benchmark & Best node & Cluster & Confidence gap & Ambiguous & Match \\
\midrule
  Integer BMU & 1 & 1 & 4+0$\sqrt{3}$ & no & $\checkmark$/$\checkmark$/$\checkmark$ \\
  Surd BMU & 6 & 3 & 10+4$\sqrt{3}$ & no & $\checkmark$/$\checkmark$/$\checkmark$ \\
  Stable tie-breaking & 1 & 7 & 0+0$\sqrt{3}$ & yes & $\checkmark$/$\checkmark$/$\checkmark$ \\
\bottomrule
\end{tabular}
}%
\end{table}



% Table 7: measured RPLU implementation artifacts
\begin{table*}[t]
\centering
\caption{Measured RPLU implementation artifacts. Metrics are reported in
each backend's native cell vocabulary; the Tang row is an arithmetic/config
probe and is not a full Pad\'{e} implementation.}
\label{tab:resource}
\resizebox{\textwidth}{!}{%
\begin{tabular}{llll}
\toprule
Artifact & Logic / registers & DSP / BRAM & Evidence \\
\midrule
Artix-7 \texttt{RPLU2PADE} & 20,277 LUTX; 6,678 FFX & 72 / 0 & routed + silicon \\
Tang \texttt{rplu2\_arith} & 9,211 LUT4; 7,926 DFF & 0 / 0 & routed + silicon \\
\midrule
\multicolumn{4}{l}{\footnotesize Artix denominator: 126,800 LUTX/FFX cells and 240 DSP48E1; timing max 36.54\,MHz, tested at 2\,MHz.} \\
\multicolumn{4}{l}{\footnotesize Tang also reports 822 ALUs; timing max 54.42/48.40\,MHz on its two reported clocks, tested at 12\,MHz.} \\
\bottomrule
\end{tabular}
}%
\end{table*}



% Table 7: Dual-Number Arithmetic — A31[epsilon]/(epsilon^2)
\begin{table}[ht]
\centering
\caption{Dual-number arithmetic over $\mathbb{A}_{\text{SPU}} =
A_{31}[\epsilon]/(\epsilon^2)$. Here $\epsilon$ is a nilpotent formal symbol
in a finite quotient ring, not an analytic infinitesimal. 9 test vectors
generated by the publication oracle; the corresponding dual-ALU RTL bench is
included in the full release gate.}
\label{tab:nsa-dual}
\resizebox{\columnwidth}{!}{%
\begin{tabular}{llllllll}
\toprule
Test & $\circ$ & Re($A$) & $\epsilon$($A$) & Re($B$) & $\epsilon$($B$) & Re($R$) & $\epsilon$($R$) \\
\midrule
  Add identity & + & (00000001,00000000,00000000,00000000) & (00000000,00000000,00000000,00000000) & (00000000,00000000,00000000,00000000) & (00000000,00000000,00000000,00000000) & (00000001,00000000,00000000,00000000) & (00000000,00000000,00000000,00000000) \\
  Add scalar (5+e7)+(3+e2) & + & (00000005,00000000,00000000,00000000) & (00000007,00000000,00000000,00000000) & (00000003,00000000,00000000,00000000) & (00000002,00000000,00000000,00000000) & (00000008,00000000,00000000,00000000) & (00000009,00000000,00000000,00000000) \\
  Add surd & + & (00000001,00000002,00000000,00000000) & (00000003,00000004,00000000,00000000) & (00000005,00000006,00000000,00000000) & (00000007,00000008,00000000,00000000) & (00000006,00000008,00000000,00000000) & (0000000A,0000000C,00000000,00000000) \\
  Mul identity & $\times$ & (00000001,00000000,00000000,00000000) & (00000000,00000000,00000000,00000000) & (00000001,00000000,00000000,00000000) & (00000000,00000000,00000000,00000000) & (00000001,00000000,00000000,00000000) & (00000000,00000000,00000000,00000000) \\
  Mul scalar (3+e0)(5+e0) & $\times$ & (00000003,00000000,00000000,00000000) & (00000000,00000000,00000000,00000000) & (00000005,00000000,00000000,00000000) & (00000000,00000000,00000000,00000000) & (0000000F,00000000,00000000,00000000) & (00000000,00000000,00000000,00000000) \\
  Mul with deriv (2+e3)(4+e5) & $\times$ & (00000002,00000000,00000000,00000000) & (00000003,00000000,00000000,00000000) & (00000004,00000000,00000000,00000000) & (00000005,00000000,00000000,00000000) & (00000008,00000000,00000000,00000000) & (00000016,00000000,00000000,00000000) \\
  Mul e$^2$=0: (0+e)(0+e)=0 & $\times$ & (00000000,00000000,00000000,00000000) & (00000001,00000000,00000000,00000000) & (00000000,00000000,00000000,00000000) & (00000001,00000000,00000000,00000000) & (00000000,00000000,00000000,00000000) & (00000000,00000000,00000000,00000000) \\
  Mul surd & $\times$ & (00000001,00000002,00000000,00000000) & (00000003,00000004,00000000,00000000) & (00000005,00000006,00000000,00000000) & (00000007,00000008,00000000,00000000) & (00000029,00000010,00000000,00000000) & (0000008E,0000003C,00000000,00000000) \\
  Mul M31 edge (P-2)(2) & $\times$ & (7FFFFFFD,00000000,00000000,00000000) & (00000000,00000000,00000000,00000000) & (00000002,00000000,00000000,00000000) & (00000000,00000000,00000000,00000000) & (7FFFFFFB,00000000,00000000,00000000) & (00000000,00000000,00000000,00000000) \\
\bottomrule
\end{tabular}
}%
\end{table}



% Table 9: RPLU v2 Pipeline Stage Latencies
\begin{table}[ht]
\centering
\caption{Contracted stage latencies for the instantiated paths.}
\label{tab:latency}
\resizebox{\columnwidth}{!}{%
\begin{tabular}{lll}
\toprule
Pipeline Stage & Module & Latency \\
\midrule
$\Phi_1$ & Kohonen SOM BMU (7-node WTA) & 5 cycles \\
$\Phi_2$ & BTU spatial router (4-lane BRAM) & 3 cycles \\
$\Phi_3$ & [4/4] Pad\'{e} Horner evaluator & 12 cycles \\
$\Phi_3$ & $A_{31}$ conjugate reduction tower & $\sim$76 cycles \\
$\Phi_4$ & Output latch & 1 cycle \\
\midrule
BTU collision queue & 64$\rightarrow$6 priority encoder & O(n) bubble stall \\
\bottomrule
\end{tabular}
}%
\end{table}



% Table 10: Reproducible verification commands
\begin{table}[ht]
\centering
\caption{Named release checks. Counts are command outputs rather than a
hand-maintained sum of selected assertions. The publication-data oracle has
39 generated checks; the project-wide gate also compiles and runs
the named RTL benches.}
\label{tab:verification}
\resizebox{\columnwidth}{!}{%
\begin{tabular}{lll}
\toprule
Command & Result & Scope \\
\midrule
\texttt{rplu\_paper\_data.py --verify} & 39/39 & table vectors \\
\texttt{spu13\_arch\_sim\_test.py} & 35/35 & adapter model \\
\texttt{test\_pade\_batch\_inversion.py} & 25/25 & tower/batch oracle \\
\texttt{test\_rational\_som.py} & 24/24 & SOM oracle \\
\texttt{run\_all\_tests.py} & 170/170 & full gate; 129 RTL \\
\bottomrule
\end{tabular}
}%
\end{table}



% ── End auto-generated data ──


% ── SECTION IV: Benchmarks ─────────────────────────────────────────

\section{Empirical Results}
\label{sec:results}

\subsection{Jet MAC Latency and Throughput}

The jet MAC pipeline (\texttt{spu13\_jet\_mac.v}) evaluates the $N=2$
Cauchy product in 12 cycles when time-division multiplexed through 16
DSP slices, or 3 cycles when fully unrolled to 96 DSP slices.
Table~\ref{tab:latency} reports cycle counts for each pipeline stage.
Under continuous operation, the TDM configuration sustains 1 MAC per
12 cycles with no pipeline bubbles (the FSM issues the next multiply
on the cycle after the previous result latches).

\subsection{Polynomial Kinematics Benchmark}

A third-order polynomial $f(t) = at^3 + bt^2 + ct + d$ is evaluated
using three jet MAC operations (chain rule through the jet algebra).
The identical computation is performed in IEEE-754 single-precision
floating-point using Horner's method:

\begin{center}
\begin{tabular}{lrr}
Metric & FP32 & Jet MAC (TDM) \\
\midrule
Latency & 18 cycles (pipelined) & 12 cycles \\
Round-trip determinism & Rounding-mode dependent & Bit-identical \\
Drift after 10$^6$ evaluations & $\pm$2--4 ULP & Zero \\
\end{tabular}
\end{center}

The jet MAC produces exact results for all rational inputs; floating-point
accumulates $\pm$2--4 ULP drift after 10$^6$ evaluations due to rounding
at each multiply-add.

\subsection{Harmonic Oscillator Propagation}

A simple harmonic oscillator $x(t) = A\cos(\omega t) + x_0$ is propagated
through 10$^3$ time steps using the jet pipeline's [4/4] Pad\'{e}
approximant as the integration kernel. The identical trajectory computed
in FP32 diverges by $\sim$0.15\% over 10$^3$ steps due to accumulated
rounding error; the exact jet trajectory converges to the closed-form
analytic solution within the rational precision of the
$A_{31}$ encoding.

\subsection{Jet-Accelerated Correction Convergence}

We compare three convergence modes for a 4-actuator correction surface
(denominator depth 16, span 75\% peak flat, material carbon/stainless):

\begin{center}
\begin{tabular}{lccc}
Strategy & Cycles & Residual & Converges? \\
\midrule
Passive SOM & 200 & 0.6400 & No \\
Active scalar ($\Delta$ only) & 11 & 0.0131 & Yes \\
Active jet ($\varepsilon^1+\varepsilon^2$) & 14 & \textbf{0.0110} & Yes \\
\end{tabular}
\end{center}

The passive SOM oscillates indefinitely---the best-matching unit drifts
across nodes without settling, producing a residual of 0.64 even after
200 cycles. Adding scalar $\Delta$-only control converges in 11 cycles to
residual 0.0131, a 98\% reduction. The full jet pipeline ($\varepsilon^1$
velocity + $\varepsilon^2$ acceleration) trades 3 additional cycles (27\%
longer) for 16\% lower residual error (0.0110 vs.\ 0.0131).

\paragraph{Why the 3-cycle tax is a win.}
In a typical discrete feedback system, higher precision requires a
multi-cycle iterative approximation loop whose cost scales with desired
accuracy. Here, the 3 additional cycles are a \emph{fixed architectural
cost}---the propagation delay as the $\varepsilon^1$ and $\varepsilon^2$
pipelines compute the damping trajectory. The jet pays this tax once,
then sustains 1 classification per 14 cycles indefinitely. The 16\%
residual reduction is the direct manifestation of the second-derivative
$\varepsilon^2$ channel damping the micro-oscillations that scalar-only
control leaves uncorrected.

\paragraph{The noise horizon hypothesis.}
While scalar-only drive converges three cycles faster under idealized,
noise-free lattice perturbations, it lacks structural awareness of
velocity and curvature. In environments with sensor noise, high-frequency
boundary jitter, or multi-node resource contention, a scalar-only approach
is structurally prone to limit-cycle oscillations---constantly
over-correcting and snapping back across the target threshold. The jet
pipeline injects a hardwired, algebraic damping term via the
$\varepsilon^2$ lane. By accounting for the instantaneous rate of change
of the error gradient, the jet core dynamically decelerates the node
weight vector as $\Delta \to 0$, acting as a local physical shock
absorber. The 16\% residual reduction measured on a clean cut is the
lower bound; under noise the gap widens.

\subsection{Lattice Collision Proximity}

Table~\ref{tab:btu-collision} shows the deterministic dispatch sequence
for single, dual, and triple simultaneous saddle-point activations.
In all cases, the dispatch order is \emph{topological} (lowest spatial
index first) and the cycle latency scales linearly with activation count.

\subsection{Pad\'{e} Rational Function Approximation}

The [4/4] Pad\'{e} evaluator computes $R(x)$ for 8 representative inputs
($x \in \{-1.0, -0.5, -0.25, 0, 0.25, 0.5, 1.0, 2.0\}$) within the fixed
~88-cycle window. RAW hazards resolve correctly when the Pad\'{e}
output is immediately consumed by a downstream QADD---the register file's
write-forwarding path delivers the result on the same cycle as writeback.

\subsection{Singular Path Absorber}

When the Pad\'{e} denominator evaluates to a zero-norm element in
$A_{31}$, the unit detector between Stage B (norm extraction) and Stage C
asserts FLAGS.V and the inversion chain is bypassed entirely.
Table~\ref{tab:singular-absorber} shows 5
scenarios: valid inversion, zero-element trap, false-positive guard
(unity), clean re-arm after exception, and post-exception valid operation.
In all cases, the hardware completes within the fixed cycle budget---no
hang, no watchdog timeout, no silent corruption.

\subsection{Montgomery Batch Inversion Speedup}

Table~\ref{tab:batch-inv} reports the cycle-model speedup of Montgomery
batch inversion vs.\ per-element tower runs, using the RTL-documented
costs of 76 cycles per tower and 3 cycles per shared multiply. The
measured workloads span single Pad\'{e} evaluations through multi-lane
manifold sweeps and trajectory propagation. At $k = 13$ (one 13-axis
manifold sweep), the batch path saves $60.0\%$ of the per-element tower
cycles at the cost of $+30.8\%$ additional MAC volume on the shared
multiplier. The MAC-volume ceiling is $\sim 33\%$, and the multiplier
is shared with the rest of the RPLU2 sidecar, so the win is real only
where the tower actually serializes---exactly the batched Pad\'{e} use
case described in Section~\ref{sec:tower}.

\begin{table}[ht]
\centering
\caption{Montgomery batch inversion cycle model: tower = 76 cycles,
shared multiply = 3 cycles. $k$ is the number of denominators inverted
in one batch. MAC $\Delta$ is the increase in shared-multiplier
operations vs.\ the per-element baseline. Values regenerated verbatim
from the \texttt{test\_pade\_batch\_inversion.py} oracle output.}
\label{tab:batch-inv}
\begin{tabular}{lrrrr}
\toprule
$k$ & Baseline cycles & Batch cycles & Speedup & MAC $\Delta$ \\
\midrule
2 (crossover)    & 206   & 139   & 1.48$\times$ & $+16.7\%$ \\
4 (quad step)    & 412   & 211   & 1.95$\times$ & $+25.0\%$ \\
13 (manifold sweep) & 1,339 & 535   & 2.50$\times$ & $+30.8\%$ \\
26 (2-step sweep)   & 2,678 & 1,003 & 2.67$\times$ & $+32.1\%$ \\
104 (trajectory)    & 10,712 & 3,811 & 2.81$\times$ & $+33.0\%$ \\
Asymptote       & ---   & ---   & 2.86$\times$ & $+33.3\%$ \\
\bottomrule
\multicolumn{5}{l}{\footnotesize Sensitivity at $k=13$: speedup is
3.20$\times$/2.50$\times$/2.12$\times$ for multiply costs of
2/3/4 cycles---the conclusion is robust to FSM handshake overhead.} \\
\end{tabular}
\end{table}

\subsection{Resource Profile}

Table~\ref{tab:resource} compares the RPLU v2 against the legacy
Morse-potential lookup table. The 3$\times$ LUT increase replaces 4
separate approximation subsystems (Morse potential, linear interpolation,
clipping, saturation) with one unified exact-rational pipeline.

Two hardware targets are characterized. On the \textbf{Xilinx Artix-7 100T}
(silicon-evidence target, 240 DSP48E1, $\sim$100\,K LUTs), the standalone
pipeline consumes $\sim$1,172 LUTs and 16 DSP slices---8\% of DSP capacity and
9\% of LUTs. Whole-image integration with live RPLU2, sidecars, and safety
layers is now treated as an Artix-7 200T / Kintex-class target because the
100T placement margin is tight for concurrent builds. On the \textbf{Gowin GW5A-25A}
(Tang Primer 25K, prototyping target), the full pipeline including rotor
core, sequencer, and VE initialization totals $\sim$5,500 LUTs (67\% of
the 8,256 LUT4s), confirming subsystem feasibility on the smaller device
but motivating larger Artix/Kintex-class devices for full concurrent operation.

\subsection{Pipeline Latency}

Table~\ref{tab:latency} details the fixed-cycle latency for each pipeline
stage. No path has variable latency---every operation completes in a
predetermined number of cycles. The longest path (Pad\'{e} denominator
inversion, $\sim$76 cycles) is still bounded and does not depend on input
values (excepting the zero-norm fast-path, which is strictly shorter).

% ── SECTION V: Discussion ──────────────────────────────────────────

\section{Discussion}
\label{sec:discussion}

\subsection{Why Not Floating-Point?}

In a standard floating-point SOM pipeline, the BMU comparison step
involves a square root ($\approx$15--30 cycles), a comparison
($\approx$1 cycle), and a potential tie-breaking heuristic (variable).
The RPLU replaces this with a single 3-stage quadrance pipeline ($\approx$5
cycles total) plus combinational WTA comparison. The difference is not
just speed---it is \emph{determinism}. An IEEE-754 square root on a value
near zero can produce a denormalized result whose rounding mode depends
on the FPU state (flush-to-zero, round-to-nearest-even, etc.). Two
identical computations on two different chips can produce different BMU
selections. The RPLU cannot.

\subsection{The Noise Horizon}

The jet pipeline's 3-cycle overhead over scalar-only control
(Section~\ref{sec:results}) is widens under noise---sensor jitter,
quantization edges, and resource contention all degrade scalar-only
convergence faster because scalar feedback lacks derivative awareness.
A scalar $\Delta$-only controller sees only the magnitude of the error,
not its rate of change; it will over-correct and oscillate across the
target boundary in proportion to the noise amplitude. The jet pipeline's
$\varepsilon^1$ (velocity) and $\varepsilon^2$ (acceleration) lanes act
as a hardwired algebraic damping term. We term this the \emph{RPLU Noise
Horizon}: the boundary beyond which scalar-only control is structurally
unstable and the jet algebra's higher-order coefficients become
necessary for deterministic convergence.

\subsection{Design Inspiration: Morphospace Analogy}

The PHSLK instruction (Section~\ref{sec:jet}) can be read as an engineering
analogy to morphogenetic closed-loop control: an Offer jet and Confirmation
jet are coherent when $O\cdot C \equiv (1,0,0)$ in
$\mathbb{A}_{\text{SPU}}$. This analogy is inspired by work on target
morphology, bioelectric networks, and anatomical homeostasis, but it is not a
formal biological claim. The hardware guarantee is narrower and purely
algebraic: the PHSLK predicate is bounded by the discrete invariants of the
local jet ring and the unit/non-unit detector in $A_{31}$.

\subsection{The Hardware Granule: Beyond Software Classification}

The confidence gap and ambiguity flag are \emph{hardware primitives}, not
post-hoc software checks. When $\text{Gap} = 0$, the hardware classifies
the result as ``boundary region'' and the output register carries both the
best and second-best cluster labels. Downstream software can consume this
ambiguity signal to trigger higher-level fusion, defer classification, or
flag the sample for expert review---without re-reading the raw feature
vector.

\section{Limitations and Future Work}
\label{sec:future}

\subsection{Dual Formal-Jet Extension}

The architecture naturally extends to the dual quotient ring
$\mathbb{A}_{\text{SPU}} = A_{31}[\epsilon]/(\epsilon^2)$, where
$\epsilon$ is a nilpotent formal generator rather than an analytic
infinitesimal. This encodes both a base value and its exact first derivative
in a single register pair. The register-pairing multiplexer
(\texttt{spu13\_nsa\_regfile\_wrapper.v}) has been designed and awaits
integration; proposed opcodes include NSA\_DQADD (0x46) and NSA\_DQMUL
(0x47) at $\sim$14--17 cycles each. For a rational function $f$ whose
required denominators are units, $f(A + \epsilon B) =
f(A) + \epsilon \cdot f'(A) \cdot B$ exactly, matching the tangent/spread
construction through finite algebra alone.

\subsection{Scoreboarded Tower Overlap}

Issue control is currently blocking during tower operations: the fixed
$\sim$88-cycle inversion window is idle time for the rest of the pipeline.
A register scoreboard permitting independent single-cycle work to issue
during tower runs would recover most of that window (e.g., evaluating
numerator Horner steps while the denominator inverts). A prototype
scoreboard module exists but requires variable-latency writeback tracking,
a same-cycle set/clear resolution rule, and a dedicated testbench before
it can be claimed; it is deliberately excluded from the verified baseline
presented here.

\subsection{Montgomery Batch Inversion RTL}

The batch inversion algorithm described in Section~\ref{sec:tower} has
been validated as a Python oracle and is under active RTL
implementation as \texttt{spu13\_batch\_inverter.v}. The module wraps the
existing conjugate reduction tower and shared M31 multiplier with a
prefix-product accumulator, an unwind FSM, and a deferred zero-divisor
isolator using the 2-multiply norm probe (stages A+B only, no Fermat).
The variable batch latency ($\sim 76 + 9(k-1)$ cycles) integrates with
the scoreboard v2's done-coupled busy bits already validated in the
baseline pipeline. This module is anticipated for the RPLU v2.1 silicon
spin on the Wukong Artix-7 100T target.

\subsection{Higher Algebraic Towers}

The conjugate reduction tower generalizes to larger finite algebras by adding
additional collapse stages. If future revisions require true fields rather
than split algebras, the basis must be rebuilt from irreducible polynomials
over M31 rather than by assuming each quadratic adjoinment remains a field.
Each additional quadratic stage adds $\sim$38 cycles to the inversion latency
and doubles the number of DSP products. For applications requiring
higher-order interactions, this preserves exact arithmetic at the cost of
linear scaling in tower height.

\subsection{Silicon Integration}

The current RTL has focused leaf-module testbenches and an actively repaired
integrated pipeline path. Two
hardware targets are in progress: a \textbf{Xilinx Artix-7 100T} (Wukong
board, 240 DSP48E1, 63,400 6-input LUTs) as the primary target for full
system concurrent operation, and a \textbf{Gowin GW5A-25A} (Tang Primer
25K, 8,256 LUT4s) as the existing component-proving board. On the 25K,
the pipeline including rotor core, sequencer, and VE initialization occupies
$\sim$5,500 LUTs (67\% of the device), confirming subsystem feasibility.
A future \textbf{RP2350B} southbridge with PSRAM will provide memory
hydration for BRAM-resident coefficient tables, eliminating the SPI flash
read bottleneck for dynamically reconfigurable material classification.

\subsection{Edge Classification and Control Primitives}

The pipeline's deterministic latency and exact arithmetic make it a natural fit
for safety-critical edge applications where floating-point rounding or timing
jitter is unacceptable. Potential directions include exact control-barrier
function kernels for real-time robotic safety filters, and geometric
classifiers for sensor fusion in aerospace guidance. The architecture is not
designed for---nor should it be evaluated against---general-purpose AI
accelerators; its value lies in the niche where reproducibility, bounded
latency, and fraction-free arithmetic are non-negotiable.

\subsection{Adjacent Application: M31 Arithmetic for Cryptographic Proof Systems}
\label{sec:zk}

Although designed for exact geometric computation, the RPLU's base-field
machinery intersects a current requirement in cryptographic proof systems.
Circle STARKs~\cite{habock2024} enable STARK provers to operate natively
over the Mersenne prime field $\mathbb{F}_{M31}$ ($p = 2^{31}-1$), chosen
for precisely the property this pipeline exploits: modular reduction by
shift-and-add within 32-bit datapaths. StarkWare's Stwo prover~\cite{stwo}
performs its soundness-critical arithmetic in QM31, a degree-4 extension of
M31, and Polygon's Plonky3 toolkit~\cite{plonky3} offers M31 among its
supported fields.

The structural alignment is at the datapath level, and one algebraic
distinction must be stated precisely. $A_{31}$ is a \emph{split} quartic
M31-algebra: because 15 is a quadratic residue modulo M31, the biquadratic
basis admits zero divisors, which is why the inversion tower carries an
explicit zero-norm detector (Section~\ref{sec:tower}). QM31, by contrast,
is a true field built from an irreducible extension tower---the property a
proof system's soundness requires. The multiplier hardware is indifferent
to this distinction: in both algebras one product is sixteen parallel
$32\times32$ base-field products with fast Mersenne reduction, followed by
a fixed linear recombination determined by the structure constants
(Section~\ref{sec:microarch}). Retargeting from $A_{31}$ to QM31 is a
rewiring of that recombination stage to the irreducible tower's constants
(rendering the zero-divisor detector vacuous), not an architectural
redesign. The RPLU pipeline above the multiplier---SOM, BTU, jet algebra,
Pad\'e evaluation---is control-loop machinery with no cryptographic role;
the reusable asset is the leaf field unit alone.

Two caveats scope the claim. First, production STARK provers demand
throughput (Circle FFT and FRI phases) that is orders of magnitude beyond a
single 16-DSP unit; this pipeline is a candidate \emph{verified building
block} for FPGA-based edge-proving experiments, not a competitive prover.
Second, the premium in this domain is correctness: proof-system soundness
depends on exact field arithmetic, so a deterministic, formally verifiable,
open-source M31 unit---with the SymbiYosys harness direction already
established for sibling modules in this repository---is the relevant
contribution, and extending bounded formal proofs to the multiplier and
inverter is prioritized future work.

% ── SECTION VI: Conclusion ─────────────────────────────────────────

\section{Conclusion}
\label{sec:conclusion}

We have presented a hardware jet algebra coprocessor over the split
biquadratic M31 algebra $A_{31}$ (Mersenne prime $p = 2^{31}-1$). The
parameterized jet MAC pipeline evaluates exact Cauchy products over the
truncated polynomial ring $A_{31}[\varepsilon]/(\varepsilon^3)$
in 12 cycles (TDM) or 3 cycles (unrolled), with bit-identical replay
across simulation runs. The architecture supports:

\begin{itemize}
\item A parameterized jet MAC computing $R_k = \sum_{i=0}^{k} J_i \cdot
K_{k-i}$ for $k \in \{0,1,2\}$, verified against 139 tests across
three independent implementation layers,
\item A 76-cycle Conjugate Reduction Tower with zero-norm singularity
detection, enabling jet inversion via a 12-cycle shadow chain,
\item A combinational BTU collision resolver with $O(n)$ priority-encoded
serial dispatch,
\item Deterministic issue control combining a write-forwarding 4R2W
register file with blocking busy/done clock-enable gating for multi-cycle
tower operations,
\item A hardware-computed Confidence Gap that elevates cluster ambiguity
from a software heuristic to a first-class architectural primitive.
\end{itemize}

The current evidence base is leaf-RTL simulation plus independent Python and
C++$\,$17 oracles; integrated place-and-route and power/timing measurement
remain future work. Two hardware targets are presented: a Xilinx Artix-7 100T
implementation target and a Gowin GW5A-25A prototype path. The pipeline is best
understood as a hardware substrate for reproducible spatial reasoning and
safety-critical control primitives; applications in edge-classification and
deterministic control-barrier kernels are natural future directions, while
general-purpose AI acceleration lies outside its scope. All RTL, software
oracles, and documentation are open-source (CC0 1.0).

\section*{Acknowledgment}
All contributions are CC0 1.0 Universal (public domain).
Open-source at: \url{github.com/pinnacletechnologysolutionsltd/SPU}.

\begin{thebibliography}{00}
\bibitem{nguyen1999}
H.~S.~Nguyen and A.~Skowron,
``Tolerance rough sets,''
in \emph{New Directions in Rough Sets, Data Mining, and Granular-Soft Computing},
Springer, 1999.

\bibitem{wheeler1945}
J.~A.~Wheeler and R.~P.~Feynman,
``Interaction with the absorber as the mechanism of radiation,''
\emph{Rev. Mod. Phys.}, vol.~17, pp.~157--181, 1945.

\bibitem{wildberger2005}
N.~J.~Wildberger,
\emph{Divine Proportions: Rational Trigonometry to Universal Geometry},
Wild Egg, 2005.

\bibitem{kohonen1995}
T.~Kohonen,
\emph{Self-Organizing Maps},
Springer, 1995.

\bibitem{pade1892}
H.~Pad\'{e},
``Sur la repr\'{e}sentation approch\'{e}e d'une fonction par des fractions rationnelles,''
\emph{Ann. Sci. \'{E}c. Norm. Sup\'{e}r.}, vol.~9, pp.~3--93, 1892.

\bibitem{montgomery1987}
P.~L.~Montgomery,
``Speeding the Pollard and elliptic curve methods of factorization,''
\emph{Math. Comp.}, vol.~48, no.~177, pp.~243--264, 1987.

\bibitem{habock2024}
U.~Hab\"{o}ck, D.~Levit, and S.~Papini,
``Circle STARKs,''
Cryptology ePrint Archive, Paper 2024/278, 2024.
\url{https://eprint.iacr.org/2024/278}

\bibitem{stwo}
StarkWare Industries,
``Stwo: a Circle-STARK prover over M31,''
\url{https://github.com/starkware-libs/stwo}

\bibitem{plonky3}
Polygon Zero,
``Plonky3: a toolkit for polynomial IOPs,''
\url{https://github.com/Plonky3/Plonky3}

\end{thebibliography}

\end{document}
